\section{\DEL{Complexity of the problem (to be removed)}\FIX{[Complexity draft --- agreed for removal]}}
\label{sec:complexity}

\begin{atodo}
You have marked this section for removal and I agree: TR-E is not a complexity
venue, the surrounding NP-hardness results in the literature you cite are
themselves conjectures, and a section that visibly ends in ``TBD'' costs more
credibility than the result would buy. The body is omitted from this build to
keep the page count down; nothing else in the paper references it.

\medskip
\textbf{Salvage one paragraph before deleting.} The charging sub-problem argument
is worth keeping as a structural remark in Section~\ref{sec:model}, because it
also motivates valid inequality \eqref{eq:vi-2}:
\end{atodo}

\begin{anew}
\noindent\textbf{Remark 2.} For a fixed charging pattern $y$ and in the absence
of Hours-of-Service constraints, an optimal solution charges at each active
station exactly enough to reach the next active station at $E^{\min}$. This
follows from the convexity and monotonicity of the cumulative charge-time
function $\mathcal T(\cdot)$: moving a unit of energy from a later station to an
earlier one moves it onto a weakly steeper part of the curve. The optimal
charging \emph{pattern} nonetheless remains combinatorial, because each active
station carries a fixed access cost $M^{stop}_i + Q_i$.
\end{anew}

\section{Solution methods} \label{sec:method}
Introducing uncertainty into a deterministic scheduling model is non-trivial: the optimal schedule depends on travel times and energy consumption that are not yet known at departure. We compare various method to tackle the problem at hand. First a simulation framework is created in Section \ref{sec:simulation} to track the real life progress of a truck, based on some decisions at each stop and realized uncertainty at each leg.
The remainder of this section is organized by \textit{when} a method commits to a decision and \textit{what} information it uses to do so, as summarized in Table \ref{tab:methodologies}: Section \ref{sec:RHLA} and Section \ref{sec:greedy} describe the two adaptive policies, which re-decide at every stop as the trip unfolds, Section \ref{sec:ro}, {Section \ref{sec:robu},} and Section \ref{sec:sp} describe the {three} anticipative methods, which fix a full schedule before departure. Finally, Section \ref{sec:oracle} introduces the perfect hindsight model that take the realized uncertainty as nominal value and solve the equivalent deterministic MILP as a benchmark for the other methods.

\begin{table} [h!]
    \centering
    \begin{tabular}{lll}\toprule
         Methodology&  Decision setting& Travel time distribution\\\midrule
         Rolling-Horizon (LA) &  Online & Scenarios\\
         Greedy&  Online & \DEL{Average}\FIX{Nominal}\\
         Robust opt. (RO)&  Offline & Robust over an interval set\\
         Budgeted robust opt. (ROBU)& Offline & Robust over a budgeted set\\
         2 stage SP&  Offline with online recourse & Scenarios\\
         Oracle&  Posteriori& Realized values\\ \bottomrule
    \end{tabular}
    \caption{Solution method characteristics}
    \label{tab:methodologies}
\end{table}

\begin{acmt}
The greedy rule uses \emph{nominal} values, not an average over a distribution:
with the feasibility guard disabled (\texttt{GUARD\_QUANTILE = None}) it uses no
distributional information at all. The reordering you already applied, so the
table now follows the subsection order, is a real readability gain.
\end{acmt}

\subsection{Simulation Framework} \label{sec:simulation}
We decouple the planning problem from the execution problem through a simulation framework, illustrated in Figure~\ref{fig:simulation_decision_selector}.

The simulation proceeds stop by stop from origin~$0$ to destination~$N$. At each stop~$i$, the vehicle has arrived with a known state $\Gamma_i = (t^a_i,\, e^a_i,\, c^d_i,\, s^d_i,\, s^w_i,\, h_i,\, \phi_i,\, n^{r2}_i, n^{\epsilon}_i)$ that summarizes accumulated time, energy, and HoS counters, where $n^{r2}_i$ and $n^{\epsilon}_i$ count the reduced rests and driving extensions already consumed. A decision method maps this state to an action $a_i$ specifying whether to charge, which break type to take, and which rest type to take. The vehicle executes $a_i$, departs, and the travel time on leg $i \to i{+}1$ is revealed as a random draw, from which the realized energy consumption follows via the ECR curve. The truck state is updated accordingly and the cycle repeats. Early arrival at a customer triggers immediate service, with the out-of-window penalty recorded\FIX{; the vehicle cannot wait for a window to open}.

\begin{figure} [h!]
    \centering
    \fcolorbox{annblue}{white}{\parbox[c][2.6cm][c]{0.9\linewidth}{\centering
    {\color{annblue}\small\textsf{FIGURE NOT IN THE REPOSITORY --- \texttt{figures/simulation.png}}}\\[3pt]
    {\color{annblue}\footnotesize Present on Overleaf; reproduced here as a placeholder only.}}}
    \caption{The simulation framework. At each stop the realized state $\Gamma_i$ is passed to a decision selector, which returns an action $a_i$ (charge / break type / rest type). The action is executed, the vehicle departs, and only then is the travel-time multiplier $\xi_i$ of the next leg drawn; the realized energy follows from the ECR curve at the implied speed. All six methods share this loop and this state representation, they differ only in the box labeled ``decision selector''.}
    \label{fig:simulation_decision_selector}
\end{figure}

\begin{acmt}
Encoding note: in the source this caption contains mojibake --- the multiplier
renders as \texttt{\^{}I\^{}$\xi$\_i} and the quotation marks as
\texttt{\^{}a\dots\^{}a}. The file has been saved once as UTF-8 and re-read as
Latin-1. Same problem in the author names (``C\^{}A\copyright line'',
``Universit\^{}A\copyright''). Re-save the \texttt{.tex} as UTF-8 in Overleaf.
\end{acmt}

The decisions methods are described below and they share the same simulation loop and vehicle state representation, only the decision step differs. The travel-time realizations are recorded throughout and fed to the oracle at the end of each run.

Because uncertainty is revealed only after departure, a leg may retroactively cause a violation that no stop-level action can repair: the consecutive-driving counter may exceed 4.5 h mid-leg or the SOC may fall below acceptable level. The simulator records these events and runs with any violation are marked infeasible.

\begin{anew}
This paragraph defines the paper's most important measurement convention in three
lines, and Section~\ref{sec:basecase} leans on it heavily. Two sentences to add:

\medskip
``Two consequences follow. First, a non-zero violation rate is unavoidable for any
online policy: a leg whose realized duration exceeds the remaining
consecutive-driving budget breaches Article 7 whatever was decided at the
preceding stop. Second, an infeasible run has no meaningful route duration and is
excluded from every gap statistic, so a method that fails often is scored only on
the realizations it survives; efficiency and reliability must therefore be read
together, Table~\ref{tab:gap} against Table~\ref{tab:feasibility}.''
\end{anew}

\subsection{Rolling-Horizon with look-Ahead policy} \label{sec:RHLA}
We adopt a \emph{rolling-horizon} policy in which a new decision is made at each stop, using only the information available at that point and considering a limited horizon to model the future. Let $i \in \{0, \ldots, N-1\}$ denote the current stop and let $\Gamma_i = (t^a_i,\, e^a_i,\, c^d_i,\, s^d_i,\, s^w_i,\, h_i,\, \phi_i,\, n^{r2}_i, n^{\epsilon}_i)$ be the vehicle state at arrival. At each stop the policy executes the following steps, which are summarized in Figure \ref{fig:RH}.

All structurally feasible actions $a \in \mathcal{A}_i$ are enumerated. An action specifies the charge binary $y \in \{0,1\}$ (relevant at CS stops only), a break type $b \in \{\varnothing, b_{45}, b_{15}, b_{30}\}$, and a rest type $r \in \{\varnothing, r_1, r_2\}$. Actions that violate structural feasibility (such as a split-break second part $b_{30}$ without a prior $b_{15}$, or a reduced rest $r_2$ when the budget is exhausted) are excluded. \NEW{No probabilistic feasibility guard is applied: an action that would strand the vehicle is not filtered a priori but is scored as infeasible in the scenarios that expose it, and is penalised through $f(a)$. This keeps the policy's information set identical to the greedy policy's.}

\begin{acmt}
The pruning sentence you commented out was right to go ---
\texttt{GUARD\_QUANTILE = None} disables that guard entirely. But the reader is
then left wondering how LA avoids stranding at all, which the purple sentence
answers.
\end{acmt}

Then, an horizon end-stop $h > i$ is determined by advancing along the route for a nominal duration of $T_{\mathrm{hor}}$ hours. This end-stop is defined via a simplistic algorithm that defines, based on rest rule, which stop can be attained in the given horizon.

For each travel leg between our current stop $i$ and the horizon end-stop $h$, a set of $S$ travel-time scenarios $\{ \boldsymbol{\xi}^{(k)}\}_{k=1}^{S}$ is drawn independently for legs $[s, h)$. Each multiplier $\xi^{(k)}_i$ is sampled from the appropriate travel time random distribution. The corresponding scenario energy $E^{(k)}_i$ is derived from the ECR curve evaluated at the implied speed.

Afterward, for each candidate action $a \in \mathcal{A}_i$ and each scenario $k$, the sub-problem MILP over the horizon $[s, h]$ is solved with $a$ fixed at the current stop and with travel times $D^{(k)}_i = D_i \xi^{(k)}_i$. The subproblem minimizes arrival time at $h$, subject to the same SOC, HoS, break, rest, and charging constraints as the full-route model. The score of action $a$ is the aggregation of the $S$ sub-problem objectives:
\[
  f(a) \;=\; \frac{1}{S}\sum_{k=1}^{S} t^{a,(k)}_h(a),
\]
with infeasible scenarios penalized by a large constant. \DEL{Alternative aggregations (worst-case, best-case, CVAR) to be considered}

\begin{acmt}
The alternative aggregations are \emph{already implemented}
(\texttt{criterion = "mean" | "worst" | "best" | "cvar"} in
\texttt{Simulation.py}, CVaR at level 0.8), so this is an un-run experiment
rather than an open question. Either report it as a one-line robustness check
alongside TB0 and MIPTAIL in Section~\ref{sec:la-config}, or delete the sentence.
\end{acmt}

Finally, the action $a^* = argmin_ {a \in \mathcal{A}_s} f(a)$ is selected, with a tie-breaking rule that favors charging ($y=1$). A final MILP re-solve on $[s, h]$ with $a^*$ fixed and nominal travel times produces the exact activity durations $(\tau^c_i, \tau^b_i, \tau^r_i)$ executed by the vehicle. The actual travel time to next stop is then drawn from the simulation framework and the vehicle state $\sigma_{i+1}$ is updated accordingly.

\begin{acmt}
$\sigma$ is already the sequential-mode binary of \eqref{eq:sigma1}. Use
$\Gamma_{i+1}$ here, as in Section~\ref{sec:simulation}. Same slip in
Section~\ref{sec:sp} (``the realized state $\sigma_i$'').
\end{acmt}

\begin{figure}[h!]
    \centering
    \fcolorbox{annblue}{white}{\parbox[c][2.6cm][c]{0.9\linewidth}{\centering
    {\color{annblue}\small\textsf{FIGURE NOT IN THE REPOSITORY --- \texttt{figures/RH.png}}}\\[3pt]
    {\color{annblue}\footnotesize Present on Overleaf; reproduced here as a placeholder only.}}}
    \caption{Rolling horizon with look ahead policy.}
    \label{fig:RH}
\end{figure}

In order to speed-up the process, only a partially relaxed LP version of the subproblems are solved. The out-of-window indicators $\delta_i$ and the next stops decisions remains as integer or binary, as well as all binary decisions at customer stops and the sequential charging binary at all CS. However, to tighten the relaxation, various valid inequalities are added to the formulation of the LP subproblems. They are summarized in \ref{sec:valid_inequ}.

\begin{acmt}
This now matches \texttt{MILP.py} exactly. Worth adding \emph{why}, in half a
sentence, because it is a real implementation finding: the big-M relaxation of
$\delta_i$ is so weak that a two-hour window miss over a hundred-hour horizon
gives $\delta \approx 0.02$, so a fully relaxed subproblem treats windows as
free and the action ranking is silently biased toward window-ignoring actions.
\end{acmt}

\subsection{Greedy heuristic} \label{sec:greedy}
To provide a baseline representative of current practice, we implement a \emph{greedy} policy that mimics the behavior of an driver operating without planning tools: stop only when forced to, charge to full whenever stopping, and avoid queuing stations known to be busy. Unlike the rolling-horizon policy, the greedy algorithm uses no look-ahead and solves no optimization sub-problem: it evaluates a fixed priority rule at each stop and returns an action immediately. A break can be inserted in the charging time if the duration is long enough.

\begin{anew}
The rule should be stated so the baseline is reproducible without reading the
code: ``The rule is evaluated in order at each stop. (1) If the shift driving
limit is about to bind, take a daily rest, reduced if the budget allows. (2)
Otherwise, if the consecutive-driving limit is about to bind, take a full break,
or the second half of a split break if the first has been taken. (3) Otherwise,
if the energy required to reach the next charging station exceeds the usable
state of charge, charge. (4) Otherwise, charge to full at any station whose
expected access delay is below a threshold. Whenever the truck is stopped at a
charging station and the charge is at least as long as the required break, the
break is declared in parallel at no additional time cost.'' Add one clause noting
that $Q_i$ is a shared model parameter, so rule (4) is not privileged
information --- otherwise a referee will suspect the baseline was handicapped.
\end{anew}

\subsection{Robust optimization} \label{sec:ro}
The deterministic MILP presented above uses nominal travel times $D_i$ on each leg $i$. In practice, travel times are uncertain and we model this via
\[
  \tilde{D}_i = D_i \,\xi_i, \qquad \xi_i \in {[\underline\xi,\ \overline\xi]}, \qquad i \in \mathcal I
\]
where $\underline\xi$ and $\overline\xi$ are fixed, physically motivated bounds on the travel-time multiplier rather than a symmetric perturbation: $\overline\xi = V^{\mathrm{nom}}/\bar v$ from a leg-average congestion floor $\bar v$, and $\underline\xi = V^{\mathrm{nom}}/\hat v$ from the vehicle speed limiter $\hat v$ (90~km/h for HGVs), so the support is asymmetric around $1$. This is the most conservative member of the budgeted family of \citet{bertsimas_price_2004}, recovered at full budget $\Gamma = N$. We retain the full box deliberately: this robust plan serves as the maximum-commitment plan. It commits every activity and every duration before departure using only the support of the uncertainty, whereas the two-stage program of Section \ref{sec:sp} exploits distributional information and online duration recourse. Section~\ref{sec:robu} relaxes this box to a budgeted uncertainty set that removes most of this conservatism while retaining a deterministic feasibility guarantee.

Because energy consumption depends on speed through the ECR curve, each $\tilde{D}_i$ also induces an uncertain energy $\tilde{E}_i(\tilde{D}_i)$. We handle this  by substituting the worst-case energy $E_i^{\max} = L_i \cdot \mathrm{ECR}\!\left(L_i / (D_i\,{\underline\xi})\right)$ into the SOC constraints, leaving them otherwise unchanged.

 The robust counterpart is then obtained directly by parameter substitution:
\[
  D_i^{\mathrm{wc}} = D_i\,{\overline\xi}, \qquad
  E_i^{\max} = L_i \cdot \mathrm{ECR}\!\left(\frac{L_i}{D_i\,{\underline\xi}}\right),
  \qquad i = 0,\ldots,N-1.
\]
The two parameters use opposite extremes of the uncertainty interval: $D_i^{\mathrm{wc}}$ uses the slowest realization, which is unfavorable for elapsed time and the Hours-of-Service accumulators, while $E_i^{\max}$ uses the fastest realization, which is unfavorable for state of charge since the ECR curve is increasing \DEL{and convex} in speed \FIX{over the operating band $[\bar v, \hat v] = [50,90]$\,km/h}. \NEW{No single realization attains both corners, so the box counterpart is conservative not only because it deviates every leg simultaneously, but because it deviates each leg in two mutually exclusive directions at once.}

\begin{acode}
$\mathrm{ECR}(v) = A/v + B + Cv^2$ is not globally increasing: with the calibrated
coefficients it has a minimum near 61\,km/h and decreases below that. It is
increasing over $[50,90]$, which is all the argument needs --- but the unqualified
claim is false and a referee checking the functional form will catch it.
\end{acode}

The resulting feasible set therefore provides a guaranteed feasibility under any realization $\xi_i \in {[\underline\xi, \overline\xi]}$.
The optimal solution at each stop is saved and used in the simulation framework to evaluate it.

\subsection{Budgeted robust optimization} \label{sec:robu}
The box counterpart of Section~\ref{sec:ro} prices every leg at its worst case simultaneously. Under independent leg multipliers this joint event has very low probability. We therefore also field the budgeted uncertainty set of \citet{bertsimas_price_2004},
\[
  \mathcal U_\Gamma = \Bigl\{ \xi \;:\; \xi_i = 1 + z_i^{+}(\overline\xi - 1) - z_i^{-}(1-\underline\xi),
  \;\; z^{+}\!, z^{-} \in [0,1]^N,
  \;\; \textstyle\sum_i z_i^{+} \le \Gamma,
  \;\; \sum_i z_i^{-} \le \Gamma \Bigr\},
\]
in which at most $\Gamma$ legs' worth of deviation mass is slow ($\xi_i \to \overline\xi$, as in Section~\ref{sec:ro}, unfavorable for elapsed time and the Hours-of-Service accumulators) and, independently, at most $\Gamma$ legs' worth is fast ($\xi_i \to \underline\xi$, unfavorable for state of charge through the ECR curve).
\DEL{We use the classical protection level $\Gamma = \sqrt{N}$ \citep{bertsimas_price_2004}.}
\FIX{We use the classical protection level $\Gamma = \lceil 1 + z_{1-\varepsilon}\sqrt{N} \rceil$ with $\varepsilon = 0.01$, for which any single constraint over $N$ independent, symmetrically distributed multipliers is violated with probability at most $\varepsilon$ \citep{bertsimas_price_2004}; this gives $\Gamma = 18$, $24$ and $31$ on the short, medium and long route classes respectively.}
The budget interpolates between the nominal model ($\Gamma = 0$) and the box of Section~\ref{sec:ro} ($\Gamma = N$).

\begin{acode}
\texttt{RObudget.py:classic\_gamma} implements
$\lceil 1 + z_{1-\varepsilon}\sqrt N\rceil$ with $\varepsilon = 0.01$, not
$\sqrt N$. The difference is large: on a short route with $N \approx 49$ legs,
$\sqrt N \approx 7$ whereas the implemented budget is 18, i.e.\ 37\,\% of all
legs. Your source already carries the correct sentence, commented out one line
above --- swap them. The realized $\Gamma$ values are now stated, because nothing
in the paper currently lets a reader size these models, and because they let the
reader check the 3.3\,\% ROBU stranding rate against the 1\,\% the budget
nominally buys (Section~\ref{sec:discussion}).
\end{acode}

Unlike the box, the budgeted counterpart does not reduce to parameter substitution. The model is a state-space formulation whose accumulators (driving since the last break, shift driving, state of charge) reset at stops selected by the binary plan, so each robust constraint involves the worst deviation over a decision-dependent window of legs. Classical constraint-wise dualization then yields products of dual variables and plan binaries and is impractical at our instance sizes \citep{poss_robust_2013}. The vehicle routing literature meets the same obstacle inside branch-and-price labeling algorithms \citep{munari_robust_2019}.
We instead solve the counterpart exactly by scenario generation \citep{mutapcic_cutting-set_2009, zeng_solving_2013}: a Master problem optimizes the plan against a finite subset $\widehat{\mathcal U} \subset \mathcal U_\Gamma$, the extensive-form model of Section~\ref{sec:sp} with a min--max objective and a single non-anticipative duration vector, each scenario carrying its own copy of the state variables, and a Pessimization step searches $\mathcal U_\Gamma$ for a realization that breaks the current plan. Because the stop sequence is fixed, once the plan is fixed every accumulator window is known, and the worst attack on each constraint family is available in closed form: fully deviate the $\Gamma$ largest legs inside that window (a vertex of $\mathcal U_\Gamma$ maximizing a nonnegative linear function). Candidate attacks are generated per continuous-driving window, per shift, per customer-window prefix, and per between-charges segment, and each is certified against the full model with the plan fixed. Violated scenarios are appended to $\widehat{\mathcal U}$ and the master is re-solved. The procedure terminates finitely \citep{zeng_solving_2013}, in practice within a handful of iterations, and the master may be solved to a relative gap without losing convergence \citep{patzold_approximate_2020}. At termination the plan is provably feasible for every realization in $\mathcal U_\Gamma$.

As for the box counterpart, the resulting plan commits every activity and every duration before departure. The optimal solution at each stop is saved and executed as is in the simulation framework to evaluate it. By the calibration of $\Gamma$, realizations outside the set may occur and render the plan infeasible in the simulation framework.

\subsection{Two-stage stochastic program}
\label{sec:sp}

We extend the deterministic MILP into a two-stage stochastic program (2-SP) with fixed recourse. First-stage decisions are the binary decisions at each stop (charge yes/no, break yes/no, rest yes/no), which are fixed once, before any uncertainty is revealed, and are therefore identical across all scenarios.
Second-stage decisions are the continuous variables at each stop (charging time, break time, and rest time), which are free to adapt to each scenario's realized travel times.

Let $k = 1, \ldots, S$ index the scenarios, each a random realization of travel time on every leg, drawn independently and with no correlation assumed between legs (the same sampling assumption used for the rolling-horizon sub-problem in Section~\ref{sec:RHLA}). Denoting the first-stage binaries $y, b, r$ and the second-stage variables $\tau^{c,(k)}, \tau^{b,(k)}, \tau^{r,(k)}$, the model solves
\[
  \min_{y,b,r \,\in\, \{0,1\}} \;\; \frac{1}{S}\sum_{k=1}^{S} \bigl(t^{a,(k)}_N(y,b,r) + \beta \sum_{i \in \mathcal C}\delta_i^{(k)}\bigr)
\]
subject to the SoC, HoS, break, rest, and charging constraints of the deterministic model holding for every scenario $k$, with nominal travel times and energies replaced by their scenario-specific values. The objective is therefore the expected arrival time across scenarios.

For evaluation, only the first-stage binary decisions are retained from the offline solve. During simulation, at each stop $i$ the remaining binaries are held fixed and the continuous activity durations are re-optimized by solving the linear program over [i, N ] with the realized state $\sigma_i$ and nominal remaining travel times. Then, only the durations at stop $i$ are executed. If this LP is infeasible, a repair step solves a restricted MILP in which binary activities may be added but not removed (for example if we need an extra CS stop), with the objective to first minimize the number of added activities and then arrival time. If the repair also fails, the run is recorded as a violation of the stochastic plan. The plan's structure is thus committed offline, its durations adapt online, and its repair frequency is itself a reported metric.

\begin{acmt}
The repair step is what makes 2SP competitive, and it is also what makes its
classification as ``offline'' questionable: it fires on \prov{99.2\,\%} of
executions (Table~\ref{tab:feasibility}). Saying so here rather than only in the
results changes how the reader interprets 2SP's low short-route gaps from the
moment they first see them.
\end{acmt}

\subsection{Perfect-Hindsight Oracle} \label{sec:oracle}
Unlike the methods above, the oracle is not a candidate policy: it is an ex-post performance bound computed once the simulation run is finished, using information no policy could have had while driving.
To assess how much is lost by the lack of advance knowledge, we compare the methodologies against a perfect-hindsight oracle: a lower bound on achievable arrival time given the travel times that actually occurred during the simulation run.

After the simulation has completed, let $\tilde{D}_0, \tilde{D}_1, \ldots, \tilde{D}_{N-1}$ be the sequence of travel times that were realized on each leg. The oracle solves the full-route MILP of Section~\ref{sec:model} with these values substituted for the nominal times $D_i$. Because the realized travel times are fixed and known when the problem is solved, this is a deterministic MILP identical in structure to the nominal model. Its optimal value $t^{a,*}_N$ is the earliest possible arrival at the destination that any policy could have achieved, had it known all travel times in advance.

The oracle is not a realizable policy: no algorithm can know $\tilde{D}_i$ before driving leg~$i$. It therefore provides a lower bound against which the optimality gap of any online policy can be measured.  The hindsight gap $\Delta$ is the relative difference of the full route duration. A small $\Delta$ indicates that look-ahead scheduling recovers most of the benefit of perfect foresight. In order to speed up the MILP solve, the rolling horizon solution is used as a first feasible solution in the perfect hindsight model.

\begin{anew}
A caveat belongs here, not only in Section~\ref{sec:basecase}, because it
qualifies every number in the paper: ``The oracle is itself a mixed-integer
program and is not always solved to proven optimality within the two-hour budget.
Where it is not, $t^{a,*}_N$ is the best incumbent found, so the reported gaps are
distances to that incumbent. On the long-route class the residual dual gap is
structural rather than a matter of solver budget --- it is driven by the packing of
the fifth mandatory rest --- so long-route gaps should be read as distances to a
schedule certified to within roughly \prov{9\,\%}, and never compared across route
classes.''
\end{anew}
